\section{问题重述与研究目标}

国际原油既是基础能源，也是高度金融化的全球大宗商品。其价格同时受到供需基本面、
库存与美元条件、市场预期以及地缘政治风险影响。地缘冲突发生后，供给中断预期、海运
受阻和风险溢价可能在数个交易日内推高油价与波动率；随后，外部冲击又会通过汇率与
进口成本进入国内成品油、工业品和居民消费价格，并最终作用于工业活动和经济增长。
因此，本题不能只建立单一预测模型，而需把``预测是否准确''、``事件期是否异常''、
``冲击如何传导''和``政策如何缓冲''分别识别后再衔接。

题面未给出``美伊战争''的精确起止日。本文将其操作性定义为：2026年2月28日开始的
中东地区军事行动及其引致的霍尔木兹海峡运输中断，并以6月17日通行恢复谅解备忘录
作为缓和阶段起点。由于2月28日为非交易日，日度事件研究从首个共同交易日
2026年3月2日开始。3月23日和4月7日的中国成品油临时调控属于国内政策事件，
不与外部战争事件合并编码。

据此，将研究任务重述为以下四个相互递进的问题：

\begin{enumerate}
  \item 以Brent现货价为国际油价主指标，在严格的时间滚动评价下预测未来
  1、3、6个月价格；同时利用日度事件研究、描述性自回归基准路径和条件波动模型，
  量化冲突不同阶段的异常价格变化与波动变化。
  \item 将第一问识别的供给、全球总需求和石油特定风险冲击作为统一输入，估计其经
  人民币原油成本、PPI、CPI和工业增加值传导至中国经济的方向、幅度、滞后和不确定性，
  并用季度实际GDP作低频验证。
  \item 对中国及六个主要石油进口国构造可比的燃油价格、CPI和工业活动面板，先量化
  中国是否表现出更强韧性，再通过缓冲变量交互和``关闭2026年临时调控''反事实解释
  政策作用及其延期成本。
  \item 利用题面``包括但不限于''的开放性，进一步研究极端油价的条件尾部概率，并在
  现行成品油调价机制约束上设计状态自适应价格平滑规则，为不同压力状态下的调价决策
  提供可检验方案。
\end{enumerate}

本文统一以2026年6月30日为观测截止日。日度层只承担事件识别和风险模拟，月度层
贯通四问，季度GDP不插值为月度数据。这样的分层使短期市场异常、历史结构冲击、宏观
响应和政策情景各自保持清楚的证据边界。
